\documentclass[10pt,a4paper]{article}
\usepackage{fontspec}
\usepackage{polyglossia}
\setmainlanguage{russian}
\setotherlanguage{english}
\setmainfont{Liberation Serif}
\usepackage[left=1.1cm,right=1.1cm,top=0.7cm,bottom=0.6cm]{geometry}
\usepackage{amsmath}
\usepackage{array}
\usepackage{multirow}
\setlength{\parindent}{0pt}
\pagestyle{empty}

\newcounter{zad}

% Поле ответа стоит СЛЕВА, номер и условие — справа от него.
% Внутри поля: зона под ответ и узкая полоска под плюс, разделённые
% вертикальной линией. Если у задачи несколько пунктов, зона ответа делится
% горизонтальными линиями на столько же блоков, а полоска под плюс остаётся
% ОДНОЙ на всю задачу — плюс ставится за задачу целиком.
%
% Базовыми считаем задачи 1--4 (проверяем, что их берут все); задачи со
% звёздочкой можно брать в любом порядке. Детям это говорится УСТНО,
% в раздаточном листке этой пометки нет.

\newlength{\otvsh}\setlength{\otvsh}{3.5cm}
\newlength{\plsh}\setlength{\plsh}{0.7cm}
\newlength{\polesh}\setlength{\polesh}{4.55cm}  % ширина всего поля целиком

% Поле НЕ имеет собственной высоты в миллиметрах: оно тянется ровно на
% высоту условия задачи — верх с верхом, низ с низом. Условие сначала
% набирается в бокс \uslbox, его высота ложится в \uslvys, и уже из неё
% делением на число блоков получается высота одного блока \blh.
\newsavebox{\uslbox}
\newlength{\uslvys}
\newlength{\blh}

% высота блока при #1 блоках: из общей высоты вычитаются линейки рамки
\newcommand{\calcblh}[1]{%
  \setlength{\blh}{\dimexpr(\uslvys-\arrayrulewidth)/#1-\arrayrulewidth\relax}}

% Полоска под плюс — ОДНА на всю задачу: плюс ставится за задачу целиком,
% поэтому правая ячейка объединена по всем блокам (multirow), а
% горизонтальные линии режут только левую колонку (\cline{1-1}).
\newcommand{\mt}[1]{\parbox[t][\blh][t]{\otvsh}{\small\bfseries #1\strut}}

\newcommand{\poleA}{\calcblh{1}\setlength{\tabcolsep}{3pt}%
  \parbox[t]{\polesh}{\vspace{0pt}%
  \begin{tabular}{|p{\otvsh}|p{\plsh}|}\hline
  \mt{} & \\ \hline\end{tabular}}}

\newcommand{\poleB}[2]{\calcblh{2}\setlength{\tabcolsep}{3pt}%
  \parbox[t]{\polesh}{\vspace{0pt}%
  \begin{tabular}{|p{\otvsh}|p{\plsh}|}\hline
  \mt{#1} & \multirow{2}{*}{} \\ \cline{1-1}
  \mt{#2} & \\ \hline\end{tabular}}}

\newcommand{\poleC}[3]{\calcblh{3}\setlength{\tabcolsep}{3pt}%
  \parbox[t]{\polesh}{\vspace{0pt}%
  \begin{tabular}{|p{\otvsh}|p{\plsh}|}\hline
  \mt{#1} & \multirow{3}{*}{} \\ \cline{1-1}
  \mt{#2} & \\ \cline{1-1}
  \mt{#3} & \\ \hline\end{tabular}}}

\newcommand{\poleD}[4]{\calcblh{4}\setlength{\tabcolsep}{3pt}%
  \parbox[t]{\polesh}{\vspace{0pt}%
  \begin{tabular}{|p{\otvsh}|p{\plsh}|}\hline
  \mt{#1} & \multirow{4}{*}{} \\ \cline{1-1}
  \mt{#2} & \\ \cline{1-1}
  \mt{#3} & \\ \cline{1-1}
  \mt{#4} & \\ \hline\end{tabular}}}

\newcommand{\shapka}[1]{%
  \noindent{\Large\bfseries Правда, ложь и перебор}%
  \hfill{\small Спецмат, 7\,И \quad·\quad 15 сентября \quad·\quad вариант #1}%
  \vspace{4pt}\hrule\vspace{6pt}
  \noindent Фамилия, имя:\ \hrulefill\rule[-4mm]{0pt}{10mm}
  \vspace{2pt}\hrule
  \setcounter{zad}{0}}

% Условие сначала набирается в бокс, чтобы узнать его точную высоту,
% и только потом рядом строится поле ровно такой же высоты.
\newlength{\uslsh}\setlength{\uslsh}{13.8cm}

\newcommand{\zad}[2]{%
  \stepcounter{zad}%
  \sbox{\uslbox}{\parbox[t]{\uslsh}{\textbf{\thezad.}\ #2}}%
  \setlength{\uslvys}{\dimexpr\ht\uslbox+\dp\uslbox\relax}%
  \vspace{1pt}\par\noindent
  \begin{tabular}{@{}p{4.6cm}@{\hspace{4mm}}p{\uslsh}@{}}
  \vspace{0pt}#1 & \vspace{0pt}\usebox{\uslbox}
  \end{tabular}\par}

\newcommand{\zadz}[3]{%
  \sbox{\uslbox}{\parbox[t]{\uslsh}{\textbf{#1$^{*}$.}\ #3}}%
  \setlength{\uslvys}{\dimexpr\ht\uslbox+\dp\uslbox\relax}%
  \vspace{1pt}\par\noindent
  \begin{tabular}{@{}p{4.6cm}@{\hspace{4mm}}p{\uslsh}@{}}
  \vspace{0pt}#2 & \vspace{0pt}\usebox{\uslbox}
  \end{tabular}\par}

\begin{document}

%======================= ВАРИАНТ A =======================
\shapka{A}

\zad{\poleB{а}{б}}{Известно, что говорящий лжёт.\newline
\textbf{(а)}~Лжец сказал: «моему дедушке самое меньшее сто лет». Сколько лет дедушке может быть самое большее?\newline
\textbf{(б)}~Волк сказал: «в этом доме не меньше десяти поросят». Сколько поросят в доме самое большее?}

\zad{\poleB{И}{П}}{На острове рыцарей и лжецов вы знакомитесь с двумя братьями, \textbf{И}ваном и \textbf{П}етром. Иван говорит: «хотя бы один из нас --- лжец». Кто из них кто?}

\zad{\poleA}{Маша сказала, что её старший брат Саша съел самое меньшее 12~конфет, а сам Саша сказал, что съел не больше 13. Оба всегда говорят правду, и Саша съел вдвое больше конфет, чем Маша. Сколько конфет съел Саша?}

\zad{\poleC{а}{б}{в}}{Рыцарь с планеты Крэбби сказал: «все трёхголовые животные на нашей планете очень умные». Верно ли каждое из трёх утверждений? Напишите три раза «да» или «нет».\newline
\textbf{(а)}~У всех умных животных на этой планете три головы.\newline
\textbf{(б)}~Если животное на этой планете глупое, то у него точно не три головы.\newline
\textbf{(в)}~Если животное не трёхголовое, то оно не умное.}

\zad{\poleA}{В кедре два дупла. Бельчонок, сидящий перед вторым дуплом, сказал две фразы:\newline
\textbf{(1)}~«в другом дупле нет орехов»; \textbf{(2)}~«хотя бы в одном дупле есть орехи».\newline
Рыжие бельчата всегда говорят правду, а серые всегда врут. Какого цвета этот бельчонок?}

\zad{\poleC{А}{С}{М}}{У \textbf{А}ни, \textbf{С}веты и \textbf{М}аши живут кошка, собака и черепаха --- у каждой своё животное. Аня сказала: «у Светы живёт собака». Света сказала: «у меня живёт собака». Маша сказала: «у Светы живёт кошка». Правду сказала только та девочка, у которой живёт собака. У кого кто живёт?}

\zad{\poleB{а}{б}}{Три лисы --- Алиса, Лариса и Инесса --- разговаривали на полянке. Лариса: «Алиса не самая хитрая». Алиса: «я хитрее Ларисы». Инесса: «Алиса хитрее меня». Известно, что самая хитрая лиса солгала, а остальные сказали правду.\newline
\textbf{(а)}~Может ли самой хитрой быть Алиса? \textbf{(б)}~Какая лиса самая хитрая?}

\zad{\poleA}{Десять островитян сидят в комнате. Первый сказал: «здесь нет ни одного рыцаря», второй: «здесь не более одного рыцаря», третий: «здесь не более двух рыцарей», и так дальше; десятый сказал: «здесь не более девяти рыцарей». Сколько рыцарей в комнате?}

\zad{\poleD{А}{Б}{С}{В}}{Вчера вечером каждая хозяйка где-то выгуливала свою питомицу. \textbf{А}лефтина ходила в парк. \textbf{Б}элла гуляла со свинкой Авдотьей. \textbf{С}ветлана была не с кошкой Муркой. Собака Спаня гуляла на стадионе. Кошка Мурка гуляла во дворе. Кто-то гулял в лесу. Кроме перечисленных в той же компании были девочка \textbf{В}ероника и черепаха Джаггернаут. Кто с кем и где был?}

\zadz{10}{\poleA}{На уроке физкультуры в шеренгу встали 25 учеников. Каждый из них либо отличник, который всегда говорит правду, либо хулиган, который всегда врёт. Отличник Влад встал на 13-е место. Все, кроме Влада, заявили: «между мной и Владом ровно 6 хулиганов». Сколько всего хулиганов в шеренге?}

\zadz{11}{\poleB{А}{Б}}{В Лесогории живут только эльфы и гномы. Гномы лгут, когда говорят про своё золото, а во всех остальных случаях говорят правду. Эльфы лгут, когда говорят про гномов, а во всех остальных случаях говорят правду. Однажды два лесогорца сказали:\newline
\emph{А:} «всё моё золото я украл у Дракона». \emph{Б:} «ты лжёшь». Кто из них эльф, а кто гном?}

\zadz{12}{\poleA}{Десять логиков пришли в кафе. Каждый заранее решил заказать себе кофе или чай, но планов остальных не знал. Официантка громко спросила: «ВСЕМ принести кофе?» --- и обошла логиков по одному, записывая ответы. Каждый громко и правдиво ответил «не знаю», «да» или «нет». Первые девять ответили «не знаю», а десятый сказал «да». Сколько логиков заказали себе кофе?}

\newpage

%======================= ВАРИАНТ B =======================
\shapka{B}

\zad{\poleB{а}{б}}{Известно, что говорящий лжёт.\newline
\textbf{(а)}~Лжец сказал: «моему дедушке самое меньшее шестьдесят лет». Сколько лет дедушке может быть самое большее?\newline
\textbf{(б)}~Волк сказал: «в этом доме не меньше семи поросят». Сколько поросят в доме самое большее?}

\zad{\poleB{И}{П}}{На острове рыцарей и лжецов вы знакомитесь с двумя братьями, \textbf{И}льёй и \textbf{П}авлом. Илья говорит: «хотя бы один из нас --- лжец». Кто из них кто?}

\zad{\poleA}{Настя сказала, что её старший брат Тимофей съел самое меньшее 14~конфет, а сам Тимофей сказал, что съел не больше 15. Оба всегда говорят правду, и Тимофей съел вдвое больше конфет, чем Настя. Сколько конфет съел Тимофей?}

\zad{\poleC{а}{б}{в}}{Рыцарь с планеты Крэбби сказал: «все шестикрылые животные на нашей планете очень быстрые». Верно ли каждое из трёх утверждений? Напишите три раза «да» или «нет».\newline
\textbf{(а)}~У всех быстрых животных на этой планете шесть крыльев.\newline
\textbf{(б)}~Если животное на этой планете медленное, то у него точно не шесть крыльев.\newline
\textbf{(в)}~Если животное не шестикрылое, то оно не быстрое.}

\zad{\poleA}{В сосне две норки. Хомячок, сидящий перед второй норкой, сказал две фразы:\newline
\textbf{(1)}~«в другой норке нет желудей»; \textbf{(2)}~«хотя бы в одной норке есть жёлуди».\newline
Полосатые хомячки всегда говорят правду, а пятнистые всегда врут. Какого цвета этот хомячок?}

\zad{\poleC{О}{Р}{Д}}{У \textbf{О}ли, \textbf{Р}иты и \textbf{Д}аши живут хомяк, попугай и ёжик --- у каждой своё животное. Оля сказала: «у Риты живёт попугай». Рита сказала: «у меня живёт попугай». Даша сказала: «у Риты живёт хомяк». Правду сказала только та девочка, у которой живёт попугай. У кого кто живёт?}

\zad{\poleB{а}{б}}{Три кошки --- Муся, Ляля и Нюша --- грелись на подоконнике. Ляля: «Муся не самая ленивая». Муся: «я ленивее Ляли». Нюша: «Муся ленивее меня». Известно, что самая ленивая кошка солгала, а остальные сказали правду.\newline
\textbf{(а)}~Может ли самой ленивой быть Муся? \textbf{(б)}~Какая кошка самая ленивая?}

\zad{\poleA}{Восемь островитян сидят в комнате. Первый сказал: «здесь нет ни одного рыцаря», второй: «здесь не более одного рыцаря», третий: «здесь не более двух рыцарей», и так дальше; восьмой сказал: «здесь не более семи рыцарей». Сколько рыцарей в комнате?}

\zad{\poleD{М}{П}{Т}{Ю}}{Вчера вечером каждая хозяйка где-то выгуливала свою питомицу. \textbf{М}арина ходила в сквер. \textbf{П}олина гуляла с хомяком Тихоном. \textbf{Т}амара была не с попугаем Гошей. Кролик Пушок гулял на катке. Попугай Гоша гулял на балконе. Кто-то гулял в поле. Кроме перечисленных в той же компании были девочка \textbf{Ю}ля и ёжик Тортила. Кто с кем и где был?}

\zadz{10}{\poleA}{На уроке физкультуры в шеренгу встали 21 ученик. Каждый из них либо отличник, который всегда говорит правду, либо хулиган, который всегда врёт. Отличник Влад встал на 11-е место. Все, кроме Влада, заявили: «между мной и Владом ровно 5 хулиганов». Сколько всего хулиганов в шеренге?}

\zadz{11}{\poleB{В}{Г}}{В Горноземье живут только феи и тролли. Тролли лгут, когда говорят про своё серебро, а во всех остальных случаях говорят правду. Феи лгут, когда говорят про троллей, а во всех остальных случаях говорят правду. Однажды два горноземца сказали:\newline
\emph{В:} «всё моё серебро я украл у Великана». \emph{Г:} «ты лжёшь». Кто из них фея, а кто тролль?}

\zadz{12}{\poleA}{Двенадцать логиков пришли в кафе. Каждый заранее решил заказать себе кофе или чай, но планов остальных не знал. Официантка громко спросила: «ВСЕМ принести кофе?» --- и обошла логиков по одному, записывая ответы. Каждый громко и правдиво ответил «не знаю», «да» или «нет». Первые одиннадцать ответили «не знаю», а двенадцатый сказал «да». Сколько логиков заказали себе кофе?}

\end{document}
